\title{LongRoPE: Extending LLM Context Window Beyond 2 Million Tokens}

\begin{abstract}
Expanding the context window in large language models (LLMs) is highly sought after. Nevertheless, several challenges—including the substantial expense of fine-tuning, limited availability of lengthy textual data, and disruptive effects from novel token positions—restrict the attainable context window length to approximately 128k tokens.

In this work, we present LongRoPE, a novel approach that, for the first time, successfully scales the context window of pre-trained LLMs to a remarkable \textbf{2048k} tokens. This expansion is accomplished with just up to 1k fine-tuning steps using training lengths up to 256k tokens, all while preserving the model's effectiveness on the initial short context lengths. Our method is grounded in three principal contributions: (i) we detect and harness two distinct non-uniformities inherent in positional interpolation, employing an efficient search procedure that yields improved fine-tuning initialization and enables up to an 8$\times$ extension even when fine-tuning is not applied; (ii) we devise a staged extension method, where we initially fine-tune a 256k-token model, subsequently applying a secondary positional interpolation to this already-extended LLM to reach the target 2048k context window; and (iii) we recalibrate LongRoPE at the 8k length scale to fully restore the model's performance on short context windows. Comprehensive evaluations on LLaMA2 and Mistral across diverse benchmarks confirm the effectiveness of our approach. Importantly, LongRoPE requires only minor alterations to the positional embeddings, thereby allowing the extended models to maintain their original architecture and benefit from existing optimizations. The source code will be made publicly available at \texttt{https://github.com/microsoft/LongRoPE}.

\section{Introduction}

Although Large Language Models (LLMs) have achieved impressive results across a range of tasks~\cite{gpt4,llama2}, their application is constrained by a finite context window size, such as the 4096 token restriction in LLaMA2~\cite{llama2}. When the input exceeds this context window, LLM performance deteriorates, as the models have not been trained to handle those additional positions. This limitation creates difficulties in critical use cases, including in-context learning scenarios that require many exemplars~\cite{Huang2023FewerIM} as well as LLM-powered agents~\cite{park2023generative,madaan2023selfrefine}.

Recent studies have demonstrated that the context window of a pre-trained LLM can be scaled up to approximately 128k through fine-tuning on extended text sequences~\cite{longlora,pi,yarn,zhang2024soaring,liu2023scaling}. However, pushing the context window even further encounters three principal challenges. Firstly, newly introduced positional indices that have not been previously trained tend to generate a large number of problematic values, which in turn cause out-of-distribution complications and hinder the convergence of fine-tuning~\cite{pi}. This challenge becomes particularly severe when enlarging the window from 4k to over $>$1000k, as this process involves introducing more than 90\% novel positions. Secondly, effective fine-tuning typically necessitates the availability of training texts that match the desired context lengths. Unfortunately, current data repositories contain few texts exceeding 1000k in length, and utilizing such lengthy texts for training entails significant computational costs, often requiring immense amounts of GPU time and resources. Thirdly, scaling up to ultra-large context windows leads to a phenomenon where attention is distributed across a vastly greater number of token positions, resulting in the attention mechanism becoming overly diffused. This dispersion ultimately undermines the model’s effectiveness when working with shorter context windows, as previously noted~\cite{pi}.

A common strategy to address the initial challenge involves interpolating RoPE positional embeddings~\cite{rope,pi}, whereby new position indices are scaled down into the original pretrained interval, as depicted in Fig.~\ref{fig:overview}. Position Interpolation (PI)~\cite{pi} performs linear interpolation of RoPE’s rotational angles using the extension ratio. In contrast, NTK~\cite{ntk,dynamicntk} recommends varying the interpolation and extrapolation rates across the different dimensions of RoPE. YaRN~\cite{yarn} divides RoPE dimensions into three groups based on frequency and assigns extrapolation, NTK, and linear interpolation strategies to each group accordingly. Nevertheless, the information entropy within positional embeddings displays a \emph{complex, non-uniform} distribution in the Transformer’s design. Current methodologies do not sufficiently exploit this intricate non-uniformity, resulting in information degradation that restricts the effective length of the context window.

Section~\ref{sec:analysis} demonstrates two primary empirical observations: \textbf{(1)} Robust positional interpolation must address two types of non-uniformities—those arising from variations across RoPE dimensions and across token positions. Reduced RoPE dimensions and tokens that appear earlier in the sequence typically require less interpolation; however, the most suitable strategies are contingent on the specific sequence length targeted for extension. 
\textbf{(2)} Incorporating these non-uniform characteristics into positional interpolation enables more effective preservation of the essential information encoded by the original RoPE, especially in critical dimensions and early token positions. This approach reduces the information loss inherent to interpolation and yields superior initializations for subsequent fine-tuning. Furthermore, this strategy supports up to an 8$\times$ sequence length increase, even when no fine-tuning is applied.

Building on these observations, we present \textbf{LongRoPE}, a robust approach that successfully extends the LLM context window to surpass 2 \emph{million} tokens. The design of LongRoPE incorporates three principal advances. Firstly, {LongRoPE} leverages multidimensional and non-uniform variations within positional interpolation to their full extent. It determines optimal rescale factors for the rotation angles of RoPE for each dimension, tailored to specific token locations. Given that the search space for these rescale factors grows exponentially with the desired window extension, {LongRoPE} adopts an evolutionary search framework that integrates two optimization strategies to enhance the efficiency of this process. An example of the optimized, rescaled RoPE resulting from this search is illustrated in Fig.~\ref{fig:overview}.

Subsequently, {LongRoPE} employs a resource-efficient, incremental expansion approach to reach a 2048k context window, circumventing the necessity for fine-tuning on extremely lengthy text sequences, which are typically rare and difficult to obtain. The process initiates by identifying an appropriate 256k length using the pre-trained LLM and performing fine-tuning at this scale. Leveraging our non-uniform positional interpolation, which supports up to an 8$\times$ increase without additional fine-tuning, we proceed to identify new RoPE rescaling factors on the LLM that has been extended and fine-tuned. This methodology ultimately allows LLaMA2 and Mistral~\cite{mistral} to operate effectively with a 2048k context window.

In order to prevent losses in performance for the original, shorter context window, {LongRoPE} further tunes the RoPE rescaling coefficients within the extended LLM. Drawing from the same approach used for enlarging the window from 256k to 2048k, we now apply our search algorithm to reduce the context length to 4k and 8k on the 256k fine-tuned LLM, thereby discouraging excessive positional interpolation. When inferring sequences shorter than 8k, we apply the rescale factors identified by our search to RoPE.

Through comprehensive experimentation with multiple LLMs on a range of long-context tasks, our approach consistently proves its efficacy. Our findings indicate that LongRoPE reliably sustains low perplexity across evaluation lengths from 4k up to 2048k, secures passkey retrieval accuracy exceeding 90\%, and matches standard benchmark performance for tasks constructed within a 4096-token context window. The LongRoPE technique is compatible with any LLMs utilizing RoPE embedding. We intend to make both our code and the LongRoPE-2048k models publicly available.

\section{Non-uniformity in Positional Interpolation}
\label{sec:analysis}

\subsection{Preliminary}

Transformer architectures depend on explicit incorporation of positional data, typically achieved through position embeddings, to capture the sequential order of input tokens. In this study, we examine RoPE~\cite{rope} position embeddings, a technique that has gained prominence in contemporary LLMs. The RoPE encoding assigned to a token located at position $n$ can be succinctly expressed as:

\begin{equation}
		\fontsize{7.8}{7.8} \selectfont
	\label{eq:rope}
	\left[ cos(n\theta_0), sin(n\theta_0), cos(n\theta_1), \cdot\cdot\cdot, cos(n\theta_{d/2-1}), sin(n\theta_{d/2-1}) \right]   
\end{equation}

Here, $d$ denotes the dimensionality of the embeddings, $n\theta_i$ specifies the rotary angle corresponding to the $n$-th token’s position, and $\theta_i=\theta^{-2i/d}$ determines the frequency for rotation. Typically, RoPE sets the base value of $\theta$ to 10000 by default.

\textbf{\emph{Context window extension ratio $s$} and positional interpolation}. The extension ratio $s$ is introduced as the quotient of the new context window length $L'$ and the original window length $L$, given by $s=\frac{L'}{L}$.

To increase the context window from $L$ to $L'$, existing positional interpolation approaches recommend adjusting the rotation frequencies $\theta_i$ downward according to the extension ratio $s$. Defining $\beta=\theta^{2/d}$ and letting $\lambda$ represent the true scaling factor corresponding to $s$, we generalize these positional interpolation techniques in the following manner:

\begin{equation}	
	\fontsize{7.5}{7.5} \selectfont
	\label{eq:unified_rope}
	\begin{aligned}
		\left[cos\left(\frac{n}{\lambda(\beta)^0}\right), sin \left(\frac{n}{\lambda(\beta)^0}\right), cos\left(\frac{n}{\lambda(\beta)^1}\right), \cdot\cdot\cdot,  sin \left(\frac{n}{\lambda(\beta)^{d/2-1}}\right) \right]   
	\end{aligned}
\end{equation}

\textbf{Linear positional interpolation (PI)}. PI~\cite{pi} proposes linearly interpolating position indices so they fit within the range used during pre-training. To achieve a target extension ratio $s$, the rotation angles assigned to every position are uniformly scaled by $\lambda=s$ for all RoPE dimensions. Nonetheless, this results in positional encodings being more densely packed together, reducing the model's capability to differentiate tokens with similar positions. As a consequence, PI typically exhibits weaker performance at larger extension ratios.

\textbf{NTK-based interpolation and extrapolation}. Works such as~\cite{ntk,dynamicntk} approach RoPE through the lens of information encoding, utilizing the framework of Neural Tangent Kernel (NTK) theory~\cite{jacot2018neural,tancik2020fourier}. To address the problem of position overcrowding in PI, they advocate for distributing the burden of interpolation more evenly across the different dimensions of RoPE. Specifically, this entails applying smaller scaling to dimensions representing higher frequencies and greater scaling to those associated with lower frequencies, leading to both positional interpolation and extrapolation, described by $\lambda=s^{i}$. The enhanced dynamic NTK~\cite{dynamicntk} modifies the extension ratio for each position in accordance with the observed sequence length. In contrast to PI, which requires model fine-tuning, NTK-oriented approaches enable the extension of context windows without such fine-tuning, although they generally achieve up to a 4$\times$ expansion ratio.

\textbf{YaRN}~\cite{yarn} presents a substantial enhancement to positional interpolation by segmenting RoPE dimensions according to frequency into three distinct groups, each utilizing a tailored interpolation method. Specifically, it applies extrapolation ($\lambda$=1) to the highest frequency dimensions, employs linear interpolation (PI) for the lowest frequency dimensions, and adopts the NTK technique for those that fall between the two extremes. A central aspect of YaRN's design is this frequency-based grouping of RoPE dimensions. However, this grouping is presently determined through manual, empirically-driven processes, which could limit optimal performance when applied to novel LLM architectures.

\subsection{Study on Non-uniform Positional Interpolation}

Drawing from the approaches of NTK and YaRN, we observe that their improvements stem from introducing non-linearity, particularly through the differential treatment of frequencies in RoPE dimensions tailored for interpolation and extrapolation tasks. Nonetheless, existing non-linear functions predominantly depend on manually crafted heuristics. This observation prompts two key questions: (1) Does the current method for positional interpolation achieve optimality? (2) Might there be undiscovered forms of non-linearity?

To address these inquiries, we employ evolution search (refer to Sec.~\ref{sec:method}) in order to identify improved non-uniform positional interpolations for LLaMA2-7B. The evolution process is steered by perplexity, evaluated on 5 randomly chosen samples from the PG19~\cite{pg19} validation dataset. Our experimental investigation yields several notable observations.

\textbf{Finding 1}: \textit{RoPE dimensions exhibit substantial non-uniformities, which are not effectively handled by current positional interpolation methods.}

We determine the optimal value of $\lambda$ for each RoPE dimension as described in Eq.~\ref{eq:unified_rope}. In Table~\ref{tbl:exp1}, we present a comparison of perplexity scores for LLaMA2-7B using various methods on the PG19 and Proof-pile~\cite{proof-pile} test datasets, all evaluated without any fine-tuning. Our approach, based on a parameter search, achieves substantial improvements, indicating that both linear (PI) and non-uniform (Dynamic-NTK and YaRN) interpolation techniques fall short of optimality. Interestingly, YaRN performs worse than PI and NTK on the PG19 dataset, failing to extend to the intended context window length when models are used without fine-tuning. For instance, when applied to an 8k context, YaRN’s perplexity rises sharply beyond the 7k token mark.

Our exploration reveals that the rescaled factors $\lambda$ in Eq.~\ref{eq:unified_rope} are no longer uniform, unlike the constant scaling factor $s$ used in PI, the NTK formulation, or the group-specific scaling of YaRN. The introduction of these non-uniform rescaling factors leads to notable gains in LLaMA2’s language modeling capability, as reflected in lower perplexity for 8k and 16k context window sizes—even without further fine-tuning. The improvement can be attributed to the enhanced preservation of the original RoPE’s structure, particularly across critical dimensions, which in turn facilitates the LLM’s ability to accurately distinguish between tokens that are close in position.

\textbf{Finding 2}: \textit{RoPE for the initial tokens in the input sequence should be extrapolated with less interpolation.} 

We propose that the RoPE applied to the initial $\hat{n}$ tokens in input sequences should utilize reduced interpolation, given that these tokens are associated with higher attention scores and therefore play a pivotal role in the attention mechanism—a finding consistent with analyses in Streaming LLM~\cite{streamingllm} and LM-Infinite~\cite{han2023lminfinite}. To test this proposition, we enlarge the context window to 8k and 16k using PI and NTK techniques, while systematically varying the number of starting tokens $\hat n$ (taking values from 0, 2, ..., 256) that are exempt from interpolation. The case $\hat n=0$ corresponds to applying standard PI and NTK across the entire sequence. Table~\ref{tbl:tokenposition} illustrates two principal findings: \textbf{(1)} excluding the first few tokens from positional interpolation enhances overall model performance, and \textbf{(2)} the optimal choice for $\hat n$ is influenced by the specific context length to which the model is being extended.

\textbf{Finding 3}: \textit{Non-uniform positional interpolation effectively extends LLM context window in both fine-tuning and non-fine-tuning settings.} 

Although our experiments demonstrate that the proposed non-uniform position interpolation considerably boosts extension performance at 8k and 16k lengths without the need for fine-tuning, extending to even longer contexts does necessitate fine-tuning. Consequently, we fine-tune LLaMA2-7B using our optimized RoPE strategy to accommodate a 64k context window (refer to the Appendix for detailed configuration). As illustrated in Table~\ref{tbl:finetune}, this approach achieves markedly better results than both PI and YaRN, both prior to and following fine-tuning of LLaMA2-7B. This superiority can be attributed to the strategic deployment of non-uniform positional interpolation, which curtails information degradation and yields improved starting parameters for fine-tuning.

\textbf{Summary}. Our research identifies two distinct non-uniformities—across RoPE dimensions and token positions. By leveraging these non-uniform aspects during positional interpolation, we attain substantial gains in the context extension capabilities of large language models.

\section{LongRoPE}

\label{sec:method}
Building upon these observations, we propose LongRoPE, which initially incorporates an effective search algorithm designed to leverage both forms of non-uniformity. Utilizing this approach, LongRoPE successfully scales the context window of LLMs to more than 2 million tokens.

\subsection{Problem Formulation}

The presence of these two forms of non-uniformity greatly expands the range of potential solutions and adds layers of difficulty to the optimization process. To manage this challenge, we conceptualize the optimization of multidimensional non-uniform position interpolation as a search problem.

Given an LLM designed to handle a context window of size $L'$ and facing lengthy input documents $\mathbf{X}$—with each sequence $\mathbf{x}\in\mathbf{X}$ exceeding $L'$ tokens—we define the original rotary angle for the $i^{th}$ dimension in RoPE at position $n$ as $\frac{n}{\beta^i}$. The resulting optimization task is expressed as:

\begin{equation}
\begin{gathered}
		\label{eq:problem}

\underset {\mathbf{x}\in\mathbf{X}; \, |\mathbf{x}|\ge L'}{ \operatorname {arg\,min} } \,  \mathcal{L}\, \left( \text{LLM} (\text{RoPE}, \mathbf{X})\right),\quad \text{where}
\\
\scalebox{0.93}{
$\underset {\begin{subarray}{l} i=0,\cdot\cdot,\frac{d}{2}-1;\\ n\in[ 0,|\mathbf{x}|);\end{subarray}}{\text{RoPE}_{(n)}}= {\left[\cdots,\cos\left(\mathbb{I}(\hat{\lambda}_{i}, \hat{n})\frac{n}{\beta^i}\right), \sin\left(\mathbb{I}(\hat{\lambda}_{i}, \hat{n})\frac{n}{\beta^i}\right),\cdots\right] }$}\\
	\text{where} \quad \mathbb{I}(\hat{\lambda}_{i},\hat{n}) = \begin{cases}
	1&\mbox{ if $n<\hat{n}$}\\
{\frac{1}{\lambda}_{i}}&\mbox{ if $n\geq\hat{n}$}
\end{cases}
	\end{gathered}
\end{equation}

In this approach, we define a collection of rescaling coefficients, $\mathbb{I}(\hat{\lambda}_{i},\hat{n})$, to address two types of non-uniformity: one associated with RoPE dimensions, represented by $\hat{\lambda}_i$, and another tied to token position, indicated by $\hat{n}$. The function $\mathbb{I}(\hat{\lambda}_{i},\hat{n})$ serves to adjust the rotation angle specifically for the $i^{th}$ RoPE dimension. Here, $\hat\lambda_{i}$ represents the respective rescaling coefficient, while $\hat{n}$ acts as a positional threshold for tokens. For token positions preceding $\hat{n}$, i.e., for the first $\hat{n}-1$ positions, the rescaling factor $\hat\lambda_{i}$ is inactive, and the original rotary angle, $\frac{n}{\beta^i}$, is retained. When the position $n$ reaches or surpasses $\hat{n}$, the adjustment factor is incorporated.

Assuming a desired context window length of $L'$, our aim is to determine the best set of rescale coefficients ($\mathbb{I}(\hat{\lambda}_{0},\hat{n})$, $\mathbb{I}(\hat{\lambda}_{1},\hat{n})$, ..., $\mathbb{I}(\hat{\lambda}_{i},\hat{n})$, ...) spanning from the $1^{st}$ to $d^{th}$ dimension of RoPE. The purpose is to ensure that the adjusted $\text{RoPE}$, when utilized within the target $\text{LLM}$, yields the lowest possible next token prediction loss $\mathcal{L}$ (corresponding to perplexity) on input samples $\mathbf{X}$ of length $L'$.

\subsection{Searching the Non-uniform Position Interpolation}

To address the problem stated in Eq.~\ref{eq:problem}, we propose a straightforward but powerful approach that seeks optimal RoPE rescale factors to take full advantage of the multidimensional, non-uniform characteristics of position embeddings.

\textbf{Search space}. We construct a comprehensive search space that captures both types of non-uniformities, as depicted in Table~\ref{tbl:searchspace}. Our method permits an independent rescale factor to be identified for each RoPE embedding dimension. To streamline the structure of the search space, the optimization is performed over $\lambda_i$ and $\hat{n}$, rather than directly over $\mathbb{I}(\hat{\lambda}_{i},\hat{n})$, with the transformation $\hat{\lambda}_{i}=1/\lambda_i$. According to Table~\ref{tbl:searchspace}, $\lambda_i$ ranges from a minimum of 1.0 (representing straightforward extrapolation) up to a maximum of $s\times1.25$ (corresponding to more aggressive interpolation than PI), incremented by 0.01. Here, $s$ denotes the desired extension ratio of the context window.

The parameter $\hat{n}$ determines how many of the earliest token positions utilize the unmodified RoPE embedding, foregoing any positional interpolation. In our experiments, $\hat{n}$ is permitted to take values from the set \{0, 1, 2, 4, 8, 12, 16, 20, 24, 28, 32, 64, 128, 256\}. Notably, setting $\hat{n}=0$ applies the optimized rescale factors to every token position.

\textbf{Evolution-based search}. The search space detailed in Table~\ref{tbl:searchspace} encompasses a vast array of positional interpolation methods, making comprehensive exploration computationally infeasible. For instance, when the extension scale is set to $s=4\times$, the total number of possible configurations reaches $400^{128/2}\times14=4\times10^{167}$. The complexity of the search space rises exponentially with higher extension ratios. To overcome these scalability issues, we adopt evolution search~\cite{guo2020single} as our search strategy and implement two additional optimization methods to significantly improve the efficiency of the process. Algorithm~\ref{alg:search} outlines the full search workflow.

\textit{Optimized initial population generation}. Rather than selecting all $P$ rescale factors at random to form the initial population, we explicitly include the three RoPE rescale factors associated with PI, NTK, and YaRN as initial individuals. The other $P$-3 members of the population are produced by applying random mutations, each with a probability of $p$, to these three specific rescale factors.

\textit{Monotonically non-decreasing constraint}. Once the initial population is established, we evaluate the perplexity for each individual using the target LLM. This involves applying the specific RoPE rescale factors to the LLM and assessing the perplexity on the input $\mathbf{X}$. The top-k individuals are then selected as parents for the evolutionary process. Nonetheless, because of the immense search space, straightforward mutation and crossover operations may often yield suboptimal candidates, resulting in many redundant perplexity evaluations. This inefficiency is especially pronounced when $L'$ is large, as each perplexity calculation requires significant inference time.

To tackle this issue, we enforce a monotonicity constraint such that the rescaled RoPE factors are non-decreasing: $\lambda_i \le \lambda_{i+1}$. Only those RoPE configurations that adhere to this condition are utilized in LLM perplexity assessments, which leads to a substantial reduction in the computational search space. More precisely, we stipulate that the values of $\lambda_i$ must strictly increase across the RoPE dimensions (with $i$ ranging from 0 to 63). This requirement for monotonicity across dimensions is grounded in NTK theory~\cite{jacot2018neural,tancik2020fourier,ntk}, which posits that dimensions associated with higher frequencies (i.e., lower $i$) benefit from smaller degrees of interpolation (reflected by smaller $\lambda_i$), while dimensions corresponding to lower frequencies (i.e., higher $i$) can tolerate greater interpolation (hence, larger $\lambda_i$).

\begin{algorithm}[t]
	\small
	\caption{The search algorithm}
	\textbf{Input:} target LLM, input samples $\mathbf{X}$, population size $P$, mutation size $N_1$, crossover size $N_2$, maximum number of iterations $\mathcal{T}$, mutation probability $p$\\

	\begin{algorithmic}[1]
		\label{alg:search}
		\STATE Top-k $\gets \phi$;	
		\STATE $\text{P}_0$ $\gets$ \textit{Initialize\_population\_with\_optimization}($P$, $\mathbf{X}$, $p$); 
		\FOR{$i = 1$ \TO $\mathcal{T}$}
		\STATE \textit{Compute\_perplexity}(LLM, $\text{P}_{i-1}$, $\mathbf{X}$);
		\STATE Top-k $\gets$ \textit{Update\_Topk}(Top-k, $\text{P}_{i-1}$);
		\STATE $\text{P}_{mutation}$ $\gets$ \textit{Mutation\_with\_mono\_constraint}(Top-k, $N_1$, $p$); 
		\STATE $\text{P}_{crossover}$ $\gets$ \textit{Crossover\_with\_mono\_constraint}(Top-k, $N_2$);
		\STATE $\text{P}_i$ $\gets$ $\text{P}_{mutation} \cup \text{P}_{crossover} \cup$ Top-k;
		\ENDFOR
		\STATE Output the member from Top-k yielding the minimum perplexity;
	\end{algorithmic}
\end{algorithm}

\textbf{8$\times$ extension without fine-tuning}. Through our evolutionary search approach, we efficiently determine non-uniform RoPE rescale factors that strategically retain essential dimensions and positions, thereby reducing information degradation due to interpolation. As shown in Fig.\ref{fig:non-ft}, this strategy enables the extension of LLaMA2's context window from 4k up to 32k tokens without the need for fine-tuning. By comparison, prior methods including PI and non-uniform NTK and YaRN exhibit substantial increases in perplexity following merely a 2$\times$ extension.

\subsection{Extending LLM Context Window to 2048K}
\label{sec:extendto2048k}

\textbf{Progressive extension to 2048k}. Here, we present our approach for expanding the context window of pre-trained LLMs from the standard 4k tokens to beyond 2048k tokens. Our results show that, using non-uniform positional interpolation, it is possible to achieve an 8$\times$ increase in context length without requiring fine-tuning. However, achieving even greater extensions—such as a 512$\times$ scale-up—necessitates fine-tuning. One approach involves selecting RoPE rescaling factors suitable for the desired 2048k context length and then performing fine-tuning. Nevertheless, this strategy encounters significant difficulties because of the extremely high computational costs involved. In addition, our empirical observations suggest that fine-tuning LLMs for such substantial extension ratios is particularly challenging (refer to Appendix).

Fortunately, LongRoPE demonstrates efficacy with both base and fine-tuned, length-extended LLMs. Accordingly, we propose a progressive and efficient approach capable of reaching the target 2048k context length using only 1k fine-tuning steps, while limiting the training sequence length to 256k.

$\Diamond$ \textit{LongRoPE search for adapting pre-trained LLMs to 256k}. Using LLaMA2 as a case study, we perform a search to expand the context window to 128k and 256k tokens. This results in an extension factor of 32$\times$ for 128k and 64$\times$ for 256k, respectively.

$\Diamond$ \textit{Fine-tuning to 256k}. Subsequently, the pre-trained LLM undergoes fine-tuning to support a context window of 256k. Initially, LLaMA2 is fine-tuned over 400 iterations employing the RoPE rescaled factors corresponding to 128k. Afterward, the RoPE rescaled factors are switched to 256k in the obtained checkpoint, followed by a further 600 fine-tuning steps. Compared to directly fine-tuning to 256k, this approach demonstrates greater efficiency.

$\Diamond$ \textit{LongRoPE-based search for extending fine-tuned extended LLM to 2048k.} In the last stage, we conduct an additional search step on the previously fine-tuned LLM set at 256k tokens. This procedure allows us to achieve an exceptionally large context window of 2048k, all without the need for extra fine-tuning. As a result, the cumulative expansion factor reaches $512\times$.

\textbf{Recovery within the shorter context window}. Upon scaling the context window to an extensive length of 2048k, a decrease in model performance is observed within the initial context window. This phenomenon is documented in the literature on positional interpolation~\cite{pi}, which shows that embedding positions into significantly higher dimensions causes those corresponding to the original context window to be constrained within a very compact space. As a result, this constraint hinders the language model's effectiveness. Notably, with a 512$\times$ extension factor, positional representations within the initial 4k context window become especially compressed.

To address this issue, an additional evolution search is conducted on the expanded LLM, specifically tuning the RoPE rescale factors for shorter context lengths such as 4k and 8k. Since shorter lengths necessitate less positional interpolation, the upper limit for the searched $\lambda$ values is lowered. At inference time, the LLM adaptively modifies the relevant RoPE rescale factors.

\section{LongRoPE}

\label{sec:method}
Building on these insights, we propose LongRoPE. This approach initially devises an effective search strategy to make optimal use of the two observed non-uniformities, subsequently leveraging this method to expand the LLM context window past 2 million tokens.

\subsection{Problem Formulation}

The presence of these two non-uniformities expands the solution space significantly and adds layers of complexity to the optimization process. To manage this challenge, we recast the multidimensional non-uniform position interpolation optimization into a search problem framework.

Consider a LLM designed for a context window of length $L'$, processing lengthy documents $\mathbf{X}$ in which every sequence $\mathbf{x}\in\mathbf{X}$ contains more tokens than $L'$. Let $\frac{n}{\beta^i}$ represent the original rotary angle associated with the $i^{th}$ dimension in the RoPE embedding at token position $n$. We can thus define the optimization task in the following way:

\begin{equation}
\begin{gathered}
		\label{eq:problem}

\underset {\mathbf{x}\in\mathbf{X}; \, |\mathbf{x}|\ge L'}{ \operatorname {arg\,min} } \,  \mathcal{L}\, \left( \text{LLM} (\text{RoPE}, \mathbf{X})\right),	\text{where \;}
\\
\scalebox{0.93}{
$\underset {\begin{subarray}{l} i=0,\cdot\cdot,\frac{d}{2}-1;\\ n\in[ 0,|\mathbf{x}|);\end{subarray}}{\text{RoPE}_(n)}= {\left[\cdot\cdot,cos\left(\mathbb{I}(\hat{\lambda}_{i}, \hat{n})\cdot\frac{n}{\beta^i}\right),  sin\left(\mathbb{I}(\hat{\lambda}_{i}, \hat{n})\cdot\frac{n}{\beta^i}\right),\cdot\cdot\right] }$}\\
	\text{where \;} \mathbb{I}(\hat{\lambda}_{i},\hat{n})=\begin{cases}
	1&\mbox{\;  $n<\hat{n}$}\\
{\frac{1}{\lambda}_{i}}&\mbox{\; $n\geq\hat{n}$}
\end{cases}
	\end{gathered}
\end{equation}

In this context, we define a group of rescale factors, denoted by $\mathbb{I}(\hat{\lambda}_{i},\hat{n})$, that are designed to address the two types of non-uniformities present. The variables $\hat{\lambda_i}$ and $\hat{n}$ correspond to the irregularities associated with the RoPE dimensions and with the token positions, respectively. In particular, $\mathbb{I}(\hat{\lambda}_{i},\hat{n})$ serves to adjust the rotation angle pertaining to the $i^{th}$ dimension in RoPE, with $\hat\lambda_{i}$ providing the dimension-specific scaling and $\hat{n}$ specifying the positional threshold. For the first $\hat{n}$–1 positions in a sequence, the scaling factor $\hat\lambda_{i}$ is not applied, and the standard RoPE rotary angle $\frac{n}{\beta^i}$ remains unaltered. Once the token index reaches $n \geq \hat{n}$, the rescale factor becomes active and modifies the rotation angle accordingly.

For a given target context window length $L'$, the aim is to determine the most effective rescale factors—denoted as ($\mathbb{I}(\hat{\lambda}_{0},\hat{n})$, $\mathbb{I}(\hat{\lambda}_{1},\hat{n})$, ..., $\mathbb{I}(\hat{\lambda}_{i},\hat{n})$, etc.)—for each RoPE dimension from the first to the $d^{th}$. This approach enables the rescaled $\text{RoPE}$ in the target $\text{LLM}$ to reach the lowest possible next-token prediction loss, $\mathcal{L}$ (which corresponds to perplexity), on input sequences $\mathbf{X}$ of length $L'$.

\subsection{Searching the Non-uniform Position Interpolation}

To address the challenge outlined in Eq.~\ref{eq:problem}, we propose a straightforward yet powerful technique that identifies the most suitable RoPE rescale factors, thereby leveraging the inherent multidimensional irregularities present in positional embeddings to their fullest extent.

\textbf{Search space}. Our method encompasses a broad search space designed to capture both forms of non-uniformity. The structure of this search space is detailed in Table~\ref{tbl:searchspace}. Concretely, our approach permits the independent determination of a unique rescale factor for each dimension within the RoPE embedding. For the sake of simplifying the design of the search space, we conduct the search over $\lambda_i$ and $\hat{n}$, rather than directly searching for $\mathbb{I}(\hat{\lambda}_{i},\hat{n})$, where $\hat{\lambda}_{i}=1/\lambda_i$. As indicated in Table~\ref{tbl:searchspace}, the search for $\lambda_i$ spans from a minimum of 1.0 (corresponding to straightforward extrapolation) up to a maximum value of $s\times1.25$ (which allows for broader interpolation than Position Interpolation (PI)), using increments of 0.01. Here, $s$ refers to the ratio by which the target context window is extended.

The parameter $\hat{n}$ specifies how many leading token positions are preserved with their original RoPE embeddings, foregoing any position interpolation. In practice, $\hat{n}$ is selected from the set \{0, 1, 2, 4, 8, 12, 16, 20, 24, 28, 32, 64, 128, 256\}. If $\hat{n}=0$, then the learned rescale factors are applied across every token position.

\textbf{Evolution-based search}. The search space detailed in Table~\ref{tbl:searchspace} encompasses a vast array of positional interpolation options, making efficient navigation particularly difficult. For instance, applying an extension of $s=4\times$ results in $400^{128/2}\times14=4\times10^{167}$ possible configurations. The size of this space grows exponentially with increased extension ratios. To manage this complexity, we adopt an evolution search strategy~\cite{guo2020single} and incorporate two additional optimization techniques that significantly enhance the efficiency of the search. The entire search process is outlined in Algorithm~\ref{alg:search}.

\textit{Optimized initial population generation}. Rather than relying solely on random initialization for the population of $P$ rescale factors, we deliberately include the three RoPE rescale factors associated with PI, NTK, and YaRN as predefined members of the initial population. For the other $P$-3 individuals, we generate them by applying random mutations to these three rescale factors, each with a mutation probability of $p$.

\textit{Monotonically non-decreasing constraint}. Once the initial population is produced, we evaluate each individual by calculating the LLM perplexity. This involves applying the designated RoPE rescale factors within the target LLM and determining the perplexity on input $\mathbf{X}$. The best-performing $k$ individuals are then selected as parents for subsequent evolutionary steps. Nonetheless, due to the enormous search space, unrefined mutation and crossover operators may frequently produce suboptimal candidates, resulting in excessive and often redundant perplexity evaluations. Such inefficiencies are amplified when $L'$ is substantial, as each perplexity computation incurs significant inference costs.

To mitigate this issue, we enforce a monotonic non-decreasing constraint on the sequence of rescaled RoPE factors such that $\lambda_i \leq \lambda_{i+1}$. Only those RoPE configurations meeting this criterion are utilized for LLM perplexity assessments, which considerably limits the scope of the search space. More concretely, we stipulate that $\lambda_i$ must grow monotonically with respect to the RoPE dimension index ($i = 0, \ldots, 63$). This imposition of monotonicity across dimensions is motivated by insights from NTK theory~\cite{jacot2018neural,tancik2020fourier,ntk}, which indicate that higher-frequency, lower-dimensional components should employ less interpolation (smaller $\lambda_i$), whereas lower-frequency, higher-dimensional ones are better suited for greater interpolation (larger $\lambda_i$).

\begin{algorithm}[t]
	\small
	\caption{The search algorithm}
	\textbf{Input:} target LLM, input samples $\mathbf{X}$,   population size $P$, mutation size $N_1$, crossover size $N_2$, max iterations $\mathcal{T}$, mutate probability $p$\\

	\begin{algorithmic}[1]
		\label{alg:search}
		\STATE Top-k $\gets \phi$;
		\STATE $\text{P}_0 \gets$ \textit{Initialize\_population\_with\_optimization}($P$, $\mathbf{X}$, $p$);
		\FOR{$i = 1$ \textbf{to} $\mathcal{T}$}
		\STATE \textit{Compute\_perplexity}(LLM, $\text{P}_{i-1}$, $\mathbf{X}$);
		\STATE Top-k $\gets$ \textit{Update\_Topk}(Top-k, $\text{P}_{i-1}$);
		\STATE $\text{P}_{mutation} \gets$ \textit{Mutation\_with\_mono\_constraint}(Top-k, $N_1$, $p$);
		\STATE $\text{P}_{crossover} \gets$ \textit{Crossover\_with\_mono\_constraint}(Top-k, $N_2$);
		\STATE $\text{P}_i \gets \text{P}_{mutation} \cup \text{P}_{crossover} \cup$ Top-k;
		\ENDFOR
		\STATE Output the member of Top-k having minimal perplexity;
	\end{algorithmic}
\end{algorithm}

\textbf{8$\times$ extension without fine-tuning}. Utilizing evolutionary search, our approach discovers optimal, non-uniform RoPE rescale factors that protect crucial positional and dimensional information, thereby reducing losses resulting from interpolation. As illustrated in Fig.\ref{fig:non-ft}, this strategy enables the expansion of LLaMA2’s context window from 4k to 32k tokens with no additional fine-tuning. By comparison, other approaches like PI, as well as non-uniform NTK and YaRN, exhibit a significant increase in perplexity beyond a 2$\times$ extension.

\subsection{Extending LLM Context Window to 2048K}
\label{sec:extendto2048k}

\textbf{Progressive extension to 2048k}. Here, we present our approach for scaling the context window of pre-trained LLMs from the conventional 4k up to beyond 2048k. As previously shown, our strategy utilizing non-uniform positional interpolation enables an 8$\times$ increase in context length without requiring any fine-tuning. When pursuing more substantial extensions (such as 512$\times$), however, fine-tuning becomes indispensable. A common approach involves determining suitable RoPE rescaling factors tailored to the target 2048k context length prior to commencing fine-tuning. This, nonetheless, introduces significant obstacles, as the process demands extremely high computational resources. Furthermore, in our empirical observations, it remains difficult to achieve effective fine-tuning of LLMs under such large extension scales (refer to Appendix for details).

Fortunately, LongRoPE demonstrates strong performance with both the base and the fine-tuned extended LLMs. Consequently, we propose a streamlined, incremental approach that reaches the desired 2048k sequence length in only 1,000 fine-tuning iterations, utilizing training sequences no longer than 256k.

$\Diamond$ \textit{LongRoPE search for adapting pre-trained LLMs to a 256k context window}. Using LLaMA2 as a case study, we perform a search aiming for context window lengths of 128k and 256k. At this phase, the model’s extension ratio reaches 32$\times$ and 64$\times$, respectively.

$\Diamond$ \textit{Fine-tuning to 256k}. Subsequently, we adapt the pre-trained LLM to extend its context window to 256k through a staged fine-tuning approach. Initially, LLaMA2 undergoes 400 fine-tuning steps with RoPE rescaling configured for 128k. Following this, the RoPE rescaling parameters are switched to support 256k, and the model is further fine-tuned for an additional 600 steps. This progressive strategy offers greater efficiency compared to directly fine-tuning the model to 256k.

$\Diamond$ \textit{LongRoPE search applied to fine-tuned extended LLMs up to 2048k}. In the last step, we carry out a subsequent search procedure on the LLM that has already been fine-tuned for a 256k input length. This process increases the maximum supported context window to a remarkable 2048k, achieved without the need for any additional fine-tuning. As a result, the model attains a context extension factor of $512\times$.

\textbf{Original context window degradation}.  When scaling up to an extensive 2048k context window, we observe diminished performance within the initial context length. This phenomenon, documented in prior work on positional interpolation~\cite{pi}, arises because embedding positions at higher dimensions in the original window are compressed into a significantly reduced area, adversely impacting model effectiveness. Specifically, at a 512$\times$ extension factor, positions inside the baseline 4k context window are especially densely packed.

To address this issue, we conduct an additional evolutionary search on the extended LLM to fine-tune RoPE rescale factors specifically for shorter context lengths such as 4k and 8k. For these cases, we constrain the upper bound of the searched $\lambda$, as shorter contexts necessitate less positional interpolation. At inference time, the LLM automatically adapts the relevant RoPE rescale factors as needed.

 \section{Experiments}

\label{sec:exp}
\subsection{Setup}
\textbf{Evaluation Tasks and models}. LongRoPE is incorporated into both LLaMA2-7B and Mistral-7B architectures, with their capabilities assessed through three dimensions: (1) the perplexity exhibited by the LLMs with extended context lengths when processing lengthy documents; (2) the Passkey retrieval task, which tests the model's proficiency in locating a straightforward passkey hidden among a large volume of unrelated content; and (3) performance on established LLM benchmarks when constrained to a 4096 token context window.

\textbf{Fine-tuning}. For the LLaMA2 model, a learning rate of 2e-5 with linear decay schedule and a total global batch size of 32 are applied. The model undergoes fine-tuning for 400 steps on the Redpajama~\cite{redpajama} dataset, which is partitioned into 128k-token segments with BOS and EOS tokens appended at the boundaries. Subsequently, utilizing the resulting checkpoint, training continues for an additional 600 steps to extend the context window to 256k tokens. Training with a 128k context length is performed on 8 A100 GPUs using the distributed training infrastructure described in~\cite{cube}, while scaling to a 256k context requires 16 A100 GPUs. For Mistral, the training uses a fixed learning rate of 1e-6 with a global batch size of 64. Both the 128k and 256k models adhere to the YaRN~\cite{yarn} training configuration, involving 400 steps on the Together Computer's Long-Data Collections~\cite{mistraldata} dataset with sequences of length 16k. All Mistral training is carried out using 4 A100 GPUs.

\textbf{Search}. When the target window size is up to 256k, we set the following parameters: $P$=64, $N_1$=$N_2$=16, $p$=0.3, and $\mathcal{T}$=40. In each iteration, we choose the best 32 candidates for mutation and crossover. Perplexity evaluation is performed on five randomly selected samples from the PG19 validation set, ensuring each sample meets at least the target context length. For window sizes exceeding 512k, the values for population, mutation, and crossover are reduced by half. In this case, perplexity is assessed on three random validation samples from Pile-Books3~\cite{pile}.

\textbf{Baselines}. To achieve a 2048k context window, we further fine-tuned models that had previously been trained on 128k and 256k context lengths. This process resulted in LongRoPE-2048k (ft=128k) and LongRoPE-2048k (ft=256k), corresponding to the LLaMA2 and Mistral models, respectively. We benchmark these four variants against leading approaches for context window extension, focusing on open-source LLMs that have undergone fine-tuning following positional interpolation through methods such as PI, NTK, and YaRN. The comparative models include Together-32k~\cite{together}, Code LLaMA~\cite{codellama}, LongLoRA-full-FT-100k~\cite{longlora}, as well as YaRN-LLaMA and YaRN-Mistral~\cite{yarn}.

\subsection{Main Results}

\label{sec:mainresult}
\textbf{Long sequence language modeling up to 256k}. Our analysis commences with a comparison to leading extended LLMs for sequence modeling evaluated at a context length of 256k. To illustrate the breadth of our approach, we utilize the Proof-pile~\cite{pg19} and PG19~\cite{pile} test sets. Perplexity is measured across different context window sizes, employing a sliding window of 256 tokens. For PG19, the evaluation encompasses all 100 documents from the test partition. For Proof-pile, consistent with YaRN~\cite{yarn}, we randomly choose 10 samples, each comprising at least 128k tokens.

Table~\ref{tbl:proof-pile} and Table~\ref{tbl:pg19} present a comparison of perplexity scores for LLaMA2 and Mistral models augmented using various interpolation techniques, evaluated on the Proof-pile and PG19 datasets, respectively. Two primary findings emerge: \textbf{(1)} the extended models consistently exhibit reduced perplexity as the evaluation length increases from 4k to 256k, demonstrating improved capacity to utilize extended contexts. \textbf{(2)} Notably, even with a context window stretched up to 16$\times$ its original size—a setting where many models typically struggle to retain strong performance at shorter ranges—our LongRoPE-2048k variants surpass leading baseline models when tested at the 256k context length.

\textbf{Long sequence language modeling beyond 2000k}. To assess performance on exceptionally lengthy documents, we utilize the Books3~\cite{pile} dataset. For the sake of evaluation efficiency, 20 books are randomly sampled, each with a length greater than 2048k, and a sliding window of 256k is applied.

As detailed in Table~\ref{tbl:books3}, LongRoPE is able to expand the context window of both LLaMA2-7B and Mistral-7B up to 2048k, while maintaining perplexity scores that are on par with, or better than, existing baselines for shorter context lengths ranging from 8k to 128k. Notably, there are significant performance variations between LLaMA2 and Mistral models at the 2048k context length. Specifically, Mistral surpasses baseline models at the shorter context ranges, yet its perplexity rises above 7 when evaluated past 256k. For LLaMA2, results are consistent with previous findings: perplexity declines as context window increases, though there are slight upticks observed at 1024k and 2048k. Additionally, for LLaMA2, utilizing LongRoPE-2048k with a fine-tuning length of 256k yields superior performance compared to using 128k, attributable to the smaller secondary extension ratio (8$\times$ vs. 16$\times$). In contrast, Mistral achieves better results with a fine-tuning window of 128k. This difference can be largely attributed to the choice of training configuration: for both the 128k and 256k fine-tuning scenarios in Mistral, the YaRN protocol is employed with a training length of 16k, which subsequently limits the extent to which Mistral’s context window can be further expanded post fine-tuning.

\textbf{Passkey retrieval}. Here, we investigate how the size of the context window affects performance in generation tasks. Adopting the synthetic passkey retrieval task introduced by~\cite{passkey}, the model is required to locate a randomly inserted passkey (a five-digit number) concealed within a lengthy document. The specifics of the prompt structure are provided in the appendix. For evaluation, we conduct 10 trials of the retrieval task, with the passkey inserted at a random, uniformly chosen position throughout the available context.

Figure~\ref{fig:passkey} presents a comparison of retrieval accuracy against baseline models. While existing methods experience a sharp decline in accuracy, reaching 0 once the token length exceeds 128k, our LongRoPE-LLaMA2-2048k (ft=256k) consistently achieves high retrieval accuracy ($\ge$90\%) across a wide range, from 4k up to 2048k, despite the difficulty of recovering a passkey among millions of tokens. Similarly, LongRoPE-Mistral-2048k (ft=128k) retains perfect accuracy (100\%) up to 1800k, with a decrease to 60\% at 2048k—a result that is consistent with the trend in Table~\ref{tbl:books3}, where perplexity modestly rises at 2048k.

\textbf{Standard benchmarks within original context window}. To assess the performance of LongRoPE-2048k models under the original context window, we employ the Hugging Face Open LLM Leaderboard~\cite{huggingfaceboard} in both zero-shot and few-shot evaluation settings. Specifically, we conduct experiments with 25-shot ARC-Challenge~\cite{allenai:arc}, 10-shot HellaSwag~\cite{zellers2019hellaswag}, 5-shot MMLU~\cite{mmlu}, and 0-shot TruthfulQA~\cite{lin2021truthfulqa}.

Table~\ref{tbl:commonsense} demonstrates that our models perform similarly to others on the original benchmark intended for shorter context windows, and actually surpass the original Mistral on TruthfulQA by an improvement of +0.5\%. Although LongRoPE-LLaMA2-2048k, which underwent fine-tuning at 256k, exhibits a marginally greater decline in performance, its results still generally fall within acceptable limits across the majority of tasks.

\subsection{Ablation Results}

\textbf{Effectiveness of the second positional interpolation}. Within our progressive extension framework, the search algorithm is employed to perform a secondary non-uniform positional interpolation on the already fine-tuned, extended LLMs. To assess this method, we carry out empirical evaluations on the fine-tuned LLaMA2-256k model. We further extend this model to 512k, 1024k, and 2048k sequence lengths using both PI and YaRN techniques. As indicated in Table~\ref{tbl:secondsearch}, our approach of non-uniform positional interpolation is able to maintain stable perplexity values throughout. By contrast, both PI and YaRN exhibit rapidly increasing perplexity as the extension ratio becomes larger.

\textbf{Effectiveness of recovery at shorter context lengths.} To address the issue of reduced performance at shorter context windows, we reoptimize the RoPE scaling parameters for LongRoPE-2048k using our search methodology. In this process, we restrict the maximum scale factors in the search space, aiming to minimize excessive interpolation specifically at 4k and 8k sequence lengths. Table~\ref{tbl:shortperformance} reports perplexity metrics for LongRoPE-LLaMA2-2048k on the Proof-pile dataset at these shorter context sizes, as well as average scores on LLM benchmarks. These findings indicate that our adjustments lead to notable gains in performance for shorter contexts.

\textbf{Analysis on the two forms of non-uniformities}. To investigate the individual contributions of the two types of non-uniformity, we conduct ablation studies. Specifically, we design two experimental settings: (i) scaling LLaMA2-7B to short sequence lengths of 16k and 32k using several approaches—PI, searching solely over RoPE dimension, and searching over both forms of non-uniformity; (ii) expanding our model fine-tuned for 256k tokens to 2048k using the same strategies. Perplexity is assessed in the absence of further fine-tuning. According to Table~\ref{tbl:searchdimension}, introducing non-uniformity in the RoPE dimension yields a notable decrease in perplexity in comparison to the linear interpolation employed by PI. Moreover, incorporating non-uniformity in token position leads to clear performance improvements at 16k and 32k, although this benefit does not extend to the 2048k case—potentially due to the challenges of such extreme sequence lengths. At these long lengths, simply retaining initial tokens without interpolation proves ineffective, and we consider further exploration of this phenomenon to be a subject for future research.

\section{Related Works}

Apart from techniques that utilize position interpolation, this section reviews relevant studies employing alternative strategies.

\textbf{Retrieval-based} techniques incorporate an external memory component for storing extended context from earlier in the sequence, coupled with retrieval mechanisms to access pertinent documents during inference~\cite{longllama,wang2023augmenting,borgeaud2022improving}. Such methods generally entail substantial architectural alterations to LLMs. By comparison, our approach introduces only minimal changes to positional embeddings, resulting in a more streamlined solution. Additionally, our method is adaptable to a broader array of long-context applications beyond retrieval, including tasks like extended document summarization and few-shot learning.

\textbf{Attention-based context window extensions}. In addition to positional embedding interpolation, several studies have extended the effective input context length of LLMs by altering the attention mechanisms within the standard context window~\cite{han2023lminfinite, streamingllm,parallelwindow}. The central approach involves designing innovative attention masks to address the challenge of increased attention computation from newly introduced positions. Such methods can be effectively combined with positional interpolation techniques, as they address orthogonal aspects of context extension.

\textbf{Fine-tuning based approaches} investigate strategies for efficiently adapting pre-trained LLMs to handle extended contexts by altering position embeddings. Prior studies, such as Code LLaMA~\cite{codellama}, LLaMA2 Long~\cite{xiong2023effective}, and ScaledRoPE~\cite{liu2023scaling}, typically set a considerably large base value for RoPE and proceed to fine-tune the model at the target context length. In contrast, our approach accommodates a wide range of target lengths, and it is capable of scaling to contexts beyond 2M tokens. Given the significant GPU demands associated with fine-tuning on long contexts (exceeding 128k tokens), newer works like LongLoRA~\cite{longlora} and PoSE~\cite{zhu2023pose} have been introduced to address computational burdens. It is worth noting that our method remains orthogonal and complementary to these efficiency-oriented fine-tuning techniques.

\section{Conclusion}

In this study, we introduce LongRoPE, a technique that significantly increases the context window of LLMs to a record-breaking 2048k tokens, all while preserving their effectiveness on standard, shorter input sequences. Our approach leverages two distinct non-uniformities present in RoPE positional embeddings, identified through an efficient evolutionary search process. This methodology confers dual advantages: it yields strong initialization for subsequent fine-tuning and allows for an 8$\times$ extension of the context window even in the absence of fine-tuning. Furthermore, we develop a staged extension protocol that utilizes LLMs fine-tuned at 256k length as intermediates, ultimately achieving a 2048k context length without requiring additional rounds of fine-tuning. Comprehensive experimental evaluations demonstrate the robustness and practicality of LongRoPE. We anticipate that our LongRoPE-2048k models will catalyze the development of novel long-context applications and stimulate new research directions in the field.

\section*{Broader Impacts} The research described in this paper is intended to contribute to the progress of Machine Learning as a discipline. While our findings could have various implications for society, we do not identify any particular effects that warrant specific mention at this time.

	\balance

\end{document}